Six-player Canadian Fish is a decentralised imperfect-information team game in which card transfers
and actions are public. We present \sestina{}, an agent developed from the FishBot policy family,
together with an exact dynamic program for the uniform distribution over deals satisfying
observer-specific game constraints. The deployed agent uses approximate inference with behavioural
weighting, a linear action score, public-history tie-breaking, a half-suit contestation feature, a
deduction-state stall detector, and guarded determinized search.

Evaluation comprises \vsevenScoredCells{} scored comparison cells over \vsevenScoredGames{} games
under a preregistered protocol. Against the registered baseline, \fcheap{}, the recorded
implementation of \sestina{} produced an observed win-rate edge of \vsevenBtwoFcheap{} percentage
points relative to parity over \vsevenBtwoFcheapN{} games, with positive estimates on both holdout
banks. The reported interval, $[\vsevenBtwoFcheapLo, \vsevenBtwoFcheapHi]$, is a nominal pooled
95\% interval formed by normal approximation from the per-bank deal-bootstrap intervals: the scored
implementation retains agent state across the deals a worker thread processes, so the
deal-independence those constructions assume, and their coverage, have not been established. Every
interval reported is labelled nominal for the same reason. The
comparison with an earlier development composite is unresolved, \vsevenBtwoComposite{} points
$[\vsevenBtwoCompositeLo, \vsevenBtwoCompositeHi]$. Cross-play and rule-variant tests retain
positive performance estimates, while the worst partner-substitution difference, \vsevenSoneMin{}
points $[\vsevenOppFiveWithholderLo, \vsevenOppFiveWithholderHi]$, does not resolve the registered
$\vsevenSoneThreshold{}$-point threshold. None of \vsevenBfourArmsWord{} tested adversarial searches
achieves a positive edge against \sestina{} at the evaluated budgets.

Four candidate modifications yield no demonstrated additional improvement under the tested
conditions. Known cross-deal state persistence and the difference between scored and
isolation-tested execution modes limit implementation-level conclusions. The results establish an
observed improvement over the specified baseline under the recorded implementation; they do not
establish coverage-validated uncertainty for that improvement, near-optimality, or an upper bound on
exploitability.
